% ---------------------------------------------------------------------------
% Section IV-A Setup and Protocol. Draft of 2026-09-11 for the mentor, who
% owns this section. Replaces the 2026-08 three-backbone draft. Every number
% comes from the 2026-09-11 export (summary.json arguments, the CSVs) or from
% the mentor's attention_backend.md; bracketed items are the mentor's to fill
% (GPU per backbone, torch version). Attention backends and transformers
% versions come from attention_backend.md and are filled. Style: no
% em-dashes, no semicolons, no prose colons, no hedges, no footnotes.
%
% What Sections IV-B and IV-C assume this section states: the selection rule
% for the tables, the paired test named once, latency as wall-clock per
% environment step, the chunk size per backbone, the decoder depth per
% backbone, and the SmolVLA two-implementation caveat.
% Peak memory was not measured in any run (no memory field in the export);
% Section V states it as a limitation. No rerun is requested for it.
% Table III (backbones) was folded into the Backbones paragraph on 2026-09-15.
% ---------------------------------------------------------------------------

\subsection{Setup and Protocol}

\textbf{Backbones.} We evaluate seven released VLA policies without
retraining. CogACT (32 Llama decoder layers), OpenVLA (32 Llama-2
layers), SpatialVLA (26 Gemma2 layers), CronusVLA (a 12-layer DiT action
decoder), UniVLA (32 Emu3 layers) and MiniVLA (24 Qwen2.5 layers) run on
SimplerEnv, and OpenVLA, UniVLA and SmolVLA (32 SmolLM2 layers) run on
LIBERO, each on the checkpoint released for that benchmark. Depth pruning
acts on the decoder stack named for each. Every backbone emits one action
per call except UniVLA, whose chunk covers five steps on WidowX and ten on
LIBERO. OpenVLA requests SDPA on SimplerEnv but defaults to eager on
LIBERO.

\textbf{Environments and episodes.} SimplerEnv WidowX has four tasks at 50
episodes each, with a step cap of 60, or 120 for the eggplant task. Google
Robot, called Fractal below, has five tasks at 50 episodes each and a cap of
80. LIBERO has four suites (Long, Goal, Object, Spatial) of ten tasks at 10
episodes each and caps of 520, 300, 280 and 220. Every
configuration of a backbone replays the same task instances and initial
states, with the episode seed fixed as 42 plus the episode index. The
LIBERO checkpoints run in their training distribution and the SimplerEnv
checkpoints under its visual matching protocol, so success measures what
a trick does to a trained policy and not how it generalizes.

\textbf{Configurations.} Each backbone and environment pair runs 14
configurations, the original policy and thirteen settings of the five
other tricks. Foveation keeps a
central disc of 0.2 or 0.5 of the image area, action repeat holds each
action for two or four steps, and depth pruning removes one, two or four
layers chosen by Block Influence, scored on a separate calibration pass
for CronusVLA and, for the other backbones, once per run on the first
test frame, whose step runs unpruned. SmolVLA's depth cells remove fixed
late layers without calibration.
% 2026-09-15 (verified in summary.json and adaptive_sparse_vla/depth_prune.py, inference_openvla_libero.py): CronusVLA depth_calibration seed 10000 on the same tasks; OpenVLA, SpatialVLA, UniVLA (and CogACT, MiniVLA) record an influence array with no seed, and the hooks ride on the first real predict call of the run, which runs unpruned, the ranking then fixed (non-ctrl mode does not reset per episode). SmolVLA: calibrated false, selected_layers [30], [28, 30], [24, 26, 28, 30] of 32.
Guarded reuse runs the strict, moderate and aggressive presets of
Section~III, and temporal fusion runs three settings, motion-entropy,
task-aware, which adds the text-to-vision attention of the preceding call
to the selector, and conservative-adaptive, which reuses fewer patches
under a shorter keyframe interval. In all, 22 backbone and environment
pairs and 47,600 episodes.
% 2026-09-15: Section III-G no longer names the three fusion settings (that sentence is commented out in the mentor's draft), so the one-clause definitions live here until the mentor restores them.
% 2026-09-13: OpenVLA Long conservative-adaptive arrived (45/100), so all 22 pairs have 14 configurations. The three reruns replaced runs at the same n, so no count changed.

\textbf{Measures.} Success is the benchmark's own criterion. Latency is
wall-clock per environment step, the episode time divided by its steps,
which has the same meaning on every harness whether a call covers one step
or a chunk. We read latency only as the change from the original within one backbone, one
environment, one software stack and one GPU class, and cells that break
this rule report success alone.
Every harness other than CronusVLA also records its per-call time, which
Sections~IV-B and~IV-C use.
Avg.\ Steps is the mean episode length. Tables~\ref{tab:simplerenv-results}
and~\ref{tab:libero-results} show, for each trick, the setting with the
highest success and, on a tie, the fewer average steps.

\textbf{Pairing and significance.} Every cell compares the same episodes
under the original policy and under the trick, its tasks pooled, and a
change is called significant when an exact two-sided McNemar test on the
episodes that flipped between failure and success gives $p < 0.05$. The
figures show point estimates. Each table cell is the best of two or three
settings, so its p-value is post-selection and read as descriptive, and
Section~IV-B reads gains over all 286 settings, where the exact test admits about eight chance passes at the flip
counts observed, about four per direction.
Sixteen of the 22 significant table drops pass a false discovery rate
correction at 0.05 over the 110 table cells, and no gain does.
Run-to-run noise is read from the episodes on which no reuse gate fired,
which run the original policy at every step. Their outcome matches the
original on SpatialVLA, UniVLA, and OpenVLA on Fractal and LIBERO on all
but at most one episode in a hundred, and differs on 2 to 17 percent of
the episodes on CogACT and CronusVLA and on MiniVLA and OpenVLA on WidowX,
so on those pairs a change of a few points is within noise and only the
paired test is read.
% 2026-09-15 (verified on the per-episode records): outcome differs from the original among zero-fire guarded-reuse episodes on CogACT Fractal 1.8%, CogACT WidowX 16.5%, CronusVLA Fractal 4.6%, CronusVLA WidowX 11.3%, MiniVLA 9.1%, OpenVLA WidowX 13.4%; OpenVLA Fractal 0, OpenVLA LIBERO 0, SpatialVLA 0 and 0.5%, UniVLA 0.2 to 0.3%. Steps of both-success zero-fire episodes are identical on 97 to 100% of episodes for the reproducing pairs and on 18 to 79% for the others. SmolVLA (reconstructed stack) 19.9%, stated in the Hardware paragraph. This replaces the earlier "all but about one episode in a hundred", which held only for UniVLA and SpatialVLA.

\textbf{Hardware and implementation.} Runs were scheduled on a shared GPU
cluster with VRAM class matched within a backbone except where noted,
and the card is recorded where known, CronusVLA and UniVLA WidowX on an
RTX 5090. Three runs, UniVLA WidowX task-aware and the CronusVLA WidowX
moderate and aggressive presets, were repeated on an RTX PRO 6000 after
failing there, and SmolVLA's guarded reuse and temporal fusion, and its
depth pruning at two layers on Long and Goal and at four on Long, Goal
and Object, ran wholly or partly under a reconstructed stack on cards of
three VRAM classes, so none of these latencies is read, and these SmolVLA
cells, whose outcome differs from the original on 10 to 30 percent of the
episodes on which no gate fired, are read through the paired test alone.
Ten cells whose gates never fired, all on UniVLA, differ from their
original row by up to 1.6 percent in latency, which we take as the
latency floor.
% 2026-09-15: per-episode implementation field 'smolvla-reconstructed-sdpa-v1' on depth_pruning2 Long (17 of 100) and Goal (26), depth_pruning4 Long (62), Goal (100), Object (100), on RTX 5090, RTX PRO 6000, A6000, L40S, 6000 Ada. Depth-1 and Spatial depth cells are all legacy.
CogACT exposes
no text-to-vision attention where fusion runs, so its task-aware setting
equals motion-entropy. Collecting attention relevance for the task-aware
setting forces an SDPA decoder into eager attention, and only OpenVLA
separates the two costs, since its LIBERO original row already runs eager
and task-aware adds the same eight percent per call as on its SDPA rows.
% 2026-09-15: the ten zero-fire cells are UniVLA Goal, Object and Spatial (all three presets) and UniVLA WidowX strict; latency deviations -1.57 to +1.41 percent. The earlier floor 0.4 to 1.4 used only the three table cells among them.

